\chapter{Le C : la matière dont tout est fait}

Nkentseu s'écrit en C++, et le C++ \emph{contient} le C. Une adresse mémoire, un
tableau, une structure, un octet : tout cela vient du C. Ce chapitre le couvre en
entier. Il suppose que vous n'avez jamais programmé, mais il ne vous fera pas
perdre de temps : chaque section dit une chose et passe à la suivante.

\section{Compiler, lier}

Un processeur ne comprend que des instructions binaires. On écrit du texte, et le
\textbf{compilateur} le traduit. Deux étapes, deux familles d'erreurs.

\begin{center}
\begin{tabular}{@{}lll@{}}
\toprule
\textbf{Étape} & \textbf{Entrée $\to$ sortie} & \textbf{Erreur typique} \\
\midrule
Compilation & un \texttt{.cpp} $\to$ un \texttt{.obj} & « je ne comprends pas cette ligne » \\
Édition de liens & les \texttt{.obj} $\to$ un \texttt{.exe} & \texttt{undefined reference to \dots} \\
\bottomrule
\end{tabular}
\end{center}

\begin{nkretenir}
Une erreur de \textbf{compilation} nomme un fichier et une ligne : la faute y
est. Une erreur d'\textbf{édition de liens} ne nomme qu'un symbole : la ligne
est correcte, c'est \emph{ailleurs} qu'un corps de fonction manque. Chercher au
mauvais endroit coûte des heures ; distinguer les deux messages coûte une
seconde.
\end{nkretenir}

\section{Premier programme}

\begin{nkcode}{Exemple pedagogique}
#include <cstdio>

int main() {
    printf("Bonjour\n");
    return 0;
}
\end{nkcode}

\texttt{\#include} fait lire un autre fichier --- celui qui déclare
\texttt{printf}. \texttt{main} est le point de départ.
\texttt{\textbackslash n} est \emph{un} caractère : le passage à la ligne.
\texttt{return 0} rend le code de sortie : \textbf{zéro veut dire succès}, à
l'inverse de l'intuition, et tous les outils suivent cette convention.

Dans Nkentseu vous écrirez \texttt{nkmain} et non \texttt{main} : le dépôt
fournit un point d'entrée portable qui devient \texttt{WinMain} sur Windows et
\texttt{android\_main} sur Android.

\section{Variables et types}

Une variable est un emplacement mémoire nommé. Le type dit combien d'octets
réserver \emph{et comment les lire} : les mêmes 32 bits donnent un entier ou un
flottant selon le type. Se tromper de type, c'est lire le bon emplacement de la
mauvaise façon.

\begin{center}
\begin{tabular}{@{}llll@{}}
\toprule
\textbf{Nkentseu} & \textbf{C} & \textbf{Octets} & \textbf{Contient} \\
\midrule
\texttt{int8} \dots \texttt{int64}   & \texttt{char} \dots \texttt{long long} & 1 à 8 & entiers signés \\
\texttt{uint8} \dots \texttt{uint64} & \texttt{unsigned \dots} & 1 à 8 & entiers $\geq 0$ \\
\texttt{float32} & \texttt{float}  & 4 & $\approx 7$ chiffres significatifs \\
\texttt{float64} & \texttt{double} & 8 & $\approx 16$ chiffres \\
\texttt{usize}   & \texttt{size\_t} & 4 ou 8 & une taille en mémoire \\
--- & \texttt{bool} & 1 & \texttt{true} / \texttt{false} \\
\bottomrule
\end{tabular}
\end{center}

\begin{nkretenir}
\texttt{int} ne fait pas la même taille partout ; \texttt{int32} si. Pour un
moteur qui vise huit plateformes, un fichier écrit sur l'une doit se relire sur
l'autre. \textbf{Écrivez les types du dépôt}, pas ceux du langage.
\end{nkretenir}

\begin{nkpiege}
\textbf{Non initialisée, une variable contient les octets qui traînaient là.} Pas
zéro. En \emph{Debug} c'est souvent un motif reconnaissable et le programme
semble marcher ; en \emph{Release} ce sont de vraies valeurs d'un calcul
précédent, et le bogue change de forme à chaque exécution. Écrivez
\texttt{int32 n = 0;}, toujours.
\end{nkpiege}

\begin{nkpiege}
\textbf{Un entier non signé ne descend pas sous zéro : il repasse par le haut.}
\texttt{uint32 n = 0; --n;} donne \texttt{4\,294\,967\,295}. D'où la boucle
infinie classique :

\begin{nkcode}{Exemple pedagogique -- NE FAIT JAMAIS DE FIN}
for (uint32 i = taille - 1; i >= 0; --i) { ... }   // i >= 0 est toujours vrai
\end{nkcode}

Et si \texttt{taille} vaut \texttt{0}, \texttt{taille - 1} vaut déjà quatre
milliards. Pour reculer : un \texttt{int32}, ou une boucle en avant.
\end{nkpiege}

\subsection{\texttt{const}, et pourquoi on le met partout}

\texttt{const} dit « ceci ne changera pas ». Le compilateur le fait respecter.

\begin{nkcode}{Exemple pedagogique}
const int32 kMax = 100;          // constante
const char *nom = "Rihen";       // pointeur vers des caracteres CONSTANTS
char *const p = tampon;          // pointeur CONSTANT vers des caracteres libres
void Lire(const NkPointer &p);   // la fonction promet de ne pas modifier p
\end{nkcode}

\begin{nkretenir}
Lisez une déclaration \textbf{de droite à gauche}. \texttt{const char *p} :
\emph{p} est un pointeur vers un \texttt{char} constant. \texttt{char *const p} :
\emph{p} est un pointeur constant vers un \texttt{char}. Le premier protège ce
qui est pointé, le second protège le pointeur.
\end{nkretenir}

\section{Opérateurs}

\begin{center}
\begin{tabular}{@{}lll@{}}
\toprule
\textbf{Famille} & \textbf{Opérateurs} & \textbf{À savoir} \\
\midrule
arithmétique & \texttt{+ - * / \%} & \texttt{\%} = reste, entiers seulement \\
comparaison  & \texttt{== != < > <= >=} & rendent un booléen \\
logique      & \texttt{\&\& || !} & évaluation \emph{paresseuse} \\
affectation  & \texttt{= += -= *= /= \%=} & \\
incrément    & \texttt{++ -{}-} & préfixé ou suffixé \\
bits         & \texttt{\& | \^{} \~{} << >>} & agissent sur les bits \\
ternaire     & \texttt{c ? a : b} & un \texttt{if} qui rend une valeur \\
\bottomrule
\end{tabular}
\end{center}

\begin{nkpiege}
\textbf{\texttt{=} affecte, \texttt{==} compare.} \texttt{if (n = 0)} affecte
zéro puis teste zéro : toujours faux, et \texttt{n} est détruit.

\textbf{La division entière tronque} : \texttt{7 / 2} vaut \texttt{3}. Pour un
résultat décimal, un opérande au moins doit l'être : \texttt{7.f / 2.f}.

\textbf{L'évaluation paresseuse est une garantie}, pas une optimisation : dans
\texttt{a \&\& b}, si \texttt{a} est faux, \texttt{b} n'est pas évalué. C'est ce
qui rend sûr \texttt{if (p != nullptr \&\& p->champ > 0)}. Inversez les termes et
vous déréférencez un pointeur nul.
\end{nkpiege}

\subsection{Les opérateurs de bits, en pratique}

Ils servent aux \textbf{drapeaux} : plusieurs oui/non dans un seul entier.

\begin{nkcode}{Exemple pedagogique}
const uint32 NK_VISIBLE  = 1u << 0;   // 0001
const uint32 NK_ACTIF    = 1u << 1;   // 0010
const uint32 NK_VERROUIL = 1u << 2;   // 0100

uint32 etat = NK_VISIBLE | NK_ACTIF;  // on POSE deux drapeaux
if (etat & NK_ACTIF)      { ... }     // on TESTE
etat &= ~NK_ACTIF;                    // on RETIRE
etat ^= NK_VISIBLE;                   // on BASCULE
\end{nkcode}

\begin{nkpiege}
\texttt{\&} n'est pas \texttt{\&\&}. \texttt{etat \& NK\_ACTIF} rend un
\emph{nombre} ; \texttt{a \&\& b} rend un \emph{booléen}. Écrire \texttt{\&\&}
entre deux drapeaux teste seulement qu'ils sont non nuls --- ce qui est vrai en
permanence.
\end{nkpiege}

\section{Contrôle de flux}

\begin{nkcode}{Exemple pedagogique}
if (a) { ... } else if (b) { ... } else { ... }

for (int32 i = 0; i < n; ++i) { ... }   // init ; test avant chaque tour ; apres
while (condition)             { ... }
do { ... } while (condition);           // au moins une fois

switch (v) {
    case NK_A: ... break;               // le `break` est OBLIGATOIRE
    case NK_B: ... break;
    default:   ... break;
}

break;      // sort de la boucle ou du switch
continue;   // passe au tour suivant
\end{nkcode}

\begin{nkpiege}
\textbf{Un \texttt{case} sans \texttt{break} tombe dans le suivant} et exécute
son code : le programme fait deux choses au lieu d'une.

\textbf{Mettez toujours les accolades}, même pour une instruction. La ligne
ajoutée six mois plus tard sous un \texttt{if} sans accolades ne fait plus partie
du \texttt{if} : l'indentation dit une chose, le compilateur en comprend une
autre, et l'\oe il ne voit rien.
\end{nkpiege}

\section{Fonctions}

Une \textbf{fonction} rend une valeur ; une \textbf{procédure} rend
\texttt{void}. Le langage ne les distingue pas, l'intention si --- et elle guide
le nom.

\begin{nkcode}{Exemple pedagogique}
int32 Somme(int32 a, int32 b) { return a + b; }   // calcule
void  Afficher(const char *s) { printf("%s\n", s); }  // agit
\end{nkcode}

\begin{nkretenir}
\textbf{Les paramètres sont copiés.} Une fonction travaille sur sa propre copie
et ne peut pas abîmer vos variables par accident. Pour qu'elle en \emph{modifie}
une, il faut lui passer son adresse.
\end{nkretenir}

\subsection{Portée et durée de vie}

\begin{center}
\begin{tabular}{@{}lll@{}}
\toprule
\textbf{Déclaration} & \textbf{Vit} & \textbf{Visible} \\
\midrule
dans un bloc & jusqu'à l'accolade fermante & dans le bloc \\
\texttt{static} dans une fonction & tout le programme & dans la fonction \\
\texttt{static} au fichier & tout le programme & dans le fichier seul \\
globale & tout le programme & partout (via \texttt{extern}) \\
\bottomrule
\end{tabular}
\end{center}

\begin{nkpiege}
\textbf{Ne rendez jamais l'adresse d'une variable locale.} Elle meurt à la sortie
de la fonction ; le pointeur rendu désigne une zone réutilisée aussitôt. Le
programme ne plante pas tout de suite --- il lit des ordures plus tard, ailleurs.
\end{nkpiege}

\section{Pointeurs}

Un pointeur contient une \textbf{adresse}. Le langage ne vérifie pas qu'elle est
valide : c'est à vous.

\begin{nkcode}{Exemple pedagogique}
int32  n  = 42;
int32 *p  = &n;    // & PREND l'adresse
int32  lu = *p;    // * SUIT l'adresse -> 42
*p = 7;            // ecrit a travers : n vaut 7

void Doubler(int32 *v) {
    if (v == nullptr) { return; }   // on verifie, toujours
    *v = *v * 2;
}
Doubler(&n);       // on passe l'ADRESSE : n vaut 14
\end{nkcode}

\begin{nkpiege}
\textbf{Trois façons dont un pointeur tue un programme}, et vous les verrez
toutes :

\begin{enumerate}
\item \textbf{nul} --- \texttt{*p} quand \texttt{p} vaut \texttt{nullptr}.
      Plantage immédiat : c'est le cas \emph{heureux}, il est reproductible ;
\item \textbf{pendant} --- il désigne une zone déjà libérée. Le programme
      continue, lit des ordures, tombe ailleurs et plus tard ;
\item \textbf{débordement} --- écrire au-delà d'un tableau écrase la variable
      d'à côté. Le symptôme apparaît dans une variable que vous n'avez pas
      touchée.
\end{enumerate}

Les deux derniers sont la raison d'être de presque tout ce que le C++ ajoute.
\end{nkpiege}

\subsection{Pointeurs de fonction}

Une fonction a une adresse, elle aussi. On peut donc la ranger dans une variable
--- c'est ainsi qu'on branche un comportement sans écrire de \texttt{switch}.

\begin{nkcode}{Idiome du depot -- une commande d'editeur}
using NkCommandeFn = void (*)(void *user);   // le TYPE d'une fonction

void Quitter(void *u) { /* ... */ }
NkCommandeFn action = &Quitter;
action(nullptr);                              // on l'appelle
\end{nkcode}

Le \texttt{void *user} est le prix d'une signature C simple : il transporte
n'importe quoi, et personne ne vérifie quoi. On ne lui passe donc \emph{que}
l'objet attendu.

\section{Tableaux et chaînes}

\begin{nkcode}{Exemple pedagogique}
int32 scores[4] = {10, 20, 30, 40};   // indices de 0 a 3 ; [4] est DEHORS
scores[0] = 15;

const char *nom = "Rihen";            // octets termines par '\0'
\end{nkcode}

Un tableau est une suite \emph{contiguë}. \texttt{scores} vaut l'adresse du
premier élément ; \texttt{scores[2]} veut dire « deux éléments plus loin ».

\begin{nkpiege}
\textbf{Un tableau ne connaît pas sa taille.} Passé à une fonction, il ne reste
que l'adresse du premier élément : la longueur est perdue. C'est pourquoi les
fonctions C prennent un couple \texttt{(pointeur, nombre)}.

\textbf{Une chaîne C se termine par \texttt{'\textbackslash 0'}.} L'oublier fait
lire jusqu'au prochain zéro rencontré, parfois très loin.
\end{nkpiege}

\section{Structures, unions, énumérations}

\begin{nkcode}{Kernel/Runtime/NKEvent/src/NKEvent/NkPointerEvent.h}
struct NkPointer {
    float32 x = 0.f;
    float32 y = 0.f;
    NkPointerPhase phase = NkPointerPhase::NK_POINTER_NONE;
    uint32 id = 0;
    bool fromTouch = false;
};
\end{nkcode}

Code réel du dépôt --- celui que vous emploierez au chapitre sur la coquille.
Notez que \textbf{chaque champ est initialisé à sa déclaration} : une
\texttt{NkPointer} neuve n'a jamais de contenu indéterminé.

\begin{nkcode}{Exemple pedagogique}
NkPointer p;
p.x = 100.f;              // POINT quand on a l'objet
NkPointer *ptr = &p;
ptr->y = 50.f;            // FLECHE quand on a un pointeur ; == (*ptr).y
\end{nkcode}

Une \texttt{union} range plusieurs champs \emph{au même endroit} : un seul est
valide à la fois. Elle sert aux événements, où la charge dépend du type --- et
c'est à vous de savoir lequel est actif.

Une \texttt{enum class} nomme un ensemble fermé de valeurs :

\begin{nkcode}{Kernel/Runtime/NKEvent/src/NKEvent/NkPointerEvent.h}
enum class NkPointerPhase : uint8 {
    NK_POINTER_NONE = 0, NK_POINTER_DOWN, NK_POINTER_MOVE,
    NK_POINTER_UP,       NK_POINTER_CANCEL
};
\end{nkcode}

\begin{nkretenir}
\texttt{enum class} plutôt que \texttt{enum} nu : les valeurs sont
\emph{cloisonnées} (on écrit \texttt{NkPointerPhase::NK\_POINTER\_UP}) et ne se
convertissent pas silencieusement en entier. Deux énumérations peuvent alors
avoir un membre du même nom sans se gêner.
\end{nkretenir}

\section{Conversions}

\begin{center}
\begin{tabular}{@{}ll@{}}
\toprule
\textbf{Écriture} & \textbf{Ce qu'elle fait} \\
\midrule
\texttt{static\_cast<T>(v)}      & conversion vérifiée à la compilation \\
\texttt{reinterpret\_cast<T>(v)} & « relis ces octets comme un T » --- dangereux \\
\texttt{const\_cast<T>(v)}       & retire un \texttt{const} --- presque toujours une faute \\
\texttt{(T)v} & la forme C : fait n'importe laquelle des trois, sans le dire \\
\bottomrule
\end{tabular}
\end{center}

\begin{nkretenir}
Employez \texttt{static\_cast}. La forme C \texttt{(T)v} est plus courte et c'est
son seul avantage : elle cache \emph{laquelle} des conversions elle applique, et
donc à quel point elle est risquée.
\end{nkretenir}

\section{Mémoire : pile et tas}

Vos variables vivaient jusqu'ici sur la \textbf{pile} : créées à l'entrée d'une
fonction, détruites à sa sortie, automatiquement. Rapide, mais petite, et la
durée de vie est imposée.

Un objet qui doit survivre à la fonction qui le crée, ou qui est trop gros, vit
sur le \textbf{tas}. On le demande, on le rend. Oublier de le rendre est une
\emph{fuite}.

\begin{nkpiege}[La règle la plus dure du dépôt]
Ni \texttt{malloc}/\texttt{free}, ni \texttt{new}/\texttt{delete} bruts. Tout
passe par NKMemory :

\begin{nkcode}{Regle du depot}
void *bloc = nkentseu::memory::NkAlloc(taille);
nkentseu::memory::NkFree(bloc);
\end{nkcode}

Et le corollaire, déjà payé plusieurs fois ici : \textbf{ce qui est alloué par un
allocateur se libère par le même}. Rendre avec le \texttt{free()} de la
bibliothèque C un tampon alloué par Nkentseu corrompt le tas --- sous Windows,
l'erreur \texttt{c0000374}, qui tombe \emph{plus tard}, dans une fonction
innocente.
\end{nkpiege}

\section{Le préprocesseur}

Il agit \emph{avant} le compilateur : il remplace du texte.

\begin{center}
\begin{tabular}{@{}ll@{}}
\toprule
\textbf{Directive} & \textbf{Effet} \\
\midrule
\texttt{\#include "x.h"} & colle le contenu de \texttt{x.h} ici \\
\texttt{\#pragma once}   & ne lire ce fichier qu'une fois \\
\texttt{\#define NK\_X 1} & remplace \texttt{NK\_X} par \texttt{1} partout \\
\texttt{\#if} / \texttt{\#endif} & compile ou non selon la plateforme \\
\bottomrule
\end{tabular}
\end{center}

\begin{nkpiege}
\textbf{Une macro n'est pas une fonction : c'est un remplacement de texte.}
\texttt{\#define DOUBLE(x) x * 2} puis \texttt{DOUBLE(1 + 1)} donne
\texttt{1 + 1 * 2}, soit \texttt{3}. Parenthésez tout :
\texttt{\#define DOUBLE(x) ((x) * 2)}. Mieux : écrivez une vraie fonction, que le
compilateur vérifie.
\end{nkpiege}

\section{Compilation séparée}

L'en-tête \texttt{.h} \textbf{déclare} ; l'implémentation \texttt{.cpp}
\textbf{définit}.

\begin{nkcode}{Exemple pedagogique}
// NkCompteur.h -- ce que les autres fichiers ont le droit de savoir
#pragma once
int32 Incrementer(int32 valeur);

// NkCompteur.cpp -- comment c'est fait
#include "NkCompteur.h"
int32 Incrementer(int32 valeur) { return valeur + 1; }
\end{nkcode}

Déclarer sans définir compile parfaitement --- et échoue à l'édition de liens.
C'est la deuxième famille d'erreurs de la première page, et vous venez d'avoir le
moyen de la provoquer.

\begin{nkpratique}
Écrivez \texttt{NkVec2.h} (une structure à deux \texttt{float32} et une fonction
\texttt{Longueur}) et son \texttt{NkVec2.cpp}.

Faites-la ensuite \textbf{échouer exprès}, deux fois : retirez un point-virgule
(erreur de compilation), puis retirez le \texttt{.cpp} du projet en gardant
l'appel (erreur d'édition de liens). Lisez les deux messages. Savoir les
distinguer d'un coup d'\oe il vous fera gagner plus de temps que n'importe quelle
autre astuce de ce cours.
\end{nkpratique}

\begin{nkbilan}
Le C tient en peu de choses : des types de taille connue, des fonctions, des
pointeurs, et une distinction entre la pile et le tas. Trois idées vous
suivront partout. \textbf{Un pointeur est une adresse} --- il ne dit ni si la
chose existe encore, ni combien il y en a. \textbf{La compilation et l'édition
de liens sont deux étapes} : un symbole déclaré et jamais défini passe la
première et échoue à la seconde, ailleurs et plus tard. Et \textbf{le
préprocesseur ne comprend rien} : il remplace du texte, ce qui explique aussi
bien son utilité que ses pièges.
\end{nkbilan}

\begin{nkquestions}
\begin{enumerate}
\item Quelle est la différence entre déclarer et définir une fonction ?
\item Que vaut un pointeur après la libération de ce qu'il désignait ?
\item Pourquoi met-on \texttt{const} sur un paramètre qu'on ne modifie pas ?
\item Qu'est-ce qui distingue la pile du tas, du point de vue de la durée de
      vie ?
\item \texttt{struct} et \texttt{union} : que partagent-ils, et qu'est-ce qui
      les sépare ?
\end{enumerate}
\end{nkquestions}

\begin{nkdemo}
Prouvez qu'un pointeur pendant n'est pas détectable en le lisant.

Écrivez un programme qui alloue un entier, garde son adresse, libère la mémoire,
puis \emph{compare} le pointeur à \texttt{NULL} et affiche le résultat. Répétez
dix fois.

\textbf{Le résultat attendu, et c'est tout l'enseignement} : il n'est jamais
\texttt{NULL}, et sa valeur ne change pas. Un pointeur libéré est
\emph{indistinguable} d'un pointeur valide par examen.

\alerte{} Ne lisez pas ce que le pointeur désigne. Le programme pourrait afficher
l'ancienne valeur sans rien signaler, et vous concluriez à tort que c'est
inoffensif.
\end{nkdemo}

\begin{nklogique}
Une fonction rend l'adresse d'une variable locale. Le programme compile, tourne,
et donne le bon résultat sur votre machine.

\begin{enumerate}
\item Expliquez pourquoi il donne pourtant un résultat juste, souvent.
\item Décrivez précisément la situation où il cesse de le donner.
\item Un collègue conclut : « ça marche, donc c'est correct ». Réfutez-le en une
      phrase, puis dites quel outil ou quelle mesure aurait tranché sans
      discussion.
\end{enumerate}
\end{nklogique}
